% !TEX root = ../projekt.tex
% =============================================================================
% STAV: hotové podľa aktuálneho kódu (bod 3 zadania, 2. časť).
% =============================================================================
\chapter{Simulátor jednotky RSS}
\label{chap:simulator}

Simulátor \texttt{rss\_sim} je napísaný v~C++17. Na zachytenej prevádzke simuluje
jednotku RSS s~najviac 128 DMA kanálmi a pre každé okno vyhodnotí záťaž kanálov.
Zámerne oddeľuje \emph{rozhodovanie} (RSS) od \emph{merania}, aby nová varianta RSS
nemohla ovplyvniť spôsob merania.

\section{Architektúra}

\begin{table}[ht]
  \centering
  \small
  \caption{Hlavné časti simulátora.}
  \label{tab:subory}
  \begin{tabular}{l p{9.5cm}}
    \toprule
    Súbor & Úloha \\
    \midrule
    \texttt{settings}    & načítanie konfiguračného súboru \\
    \texttt{parser}      & prevod riadku CSV na tuple a počty paketov \\
    \texttt{flow\_expand} & rekonštrukcia prúdu paketov v~čase (halda) \\
    \texttt{rss}         & reťaz RSS: symetria, hash, RETA alebo modulo, orákulum JSQ \\
    \texttt{simulation}  & rozhodnutia tokov, počítanie paketov, okná, zápis výsledkov \\
    \texttt{metrics}     & preťaženie, Jainov index, $\chi^2$ \\
    \texttt{wrappers/}   & jednotné rozhranie ôsmich hashovacích funkcií \\
    \texttt{analyze.py}  & súhrnná tabuľka a graf \\
    \bottomrule
  \end{tabular}
\end{table}

Beh sa nastavuje jediným súborom:

\begin{lstlisting}
data      = data/odid1.csv
out       = results/result.csv
hash      = toeplitz        # jhash chaskey halfsiphash simpletab ...
tuple     = full            # ip
symmetry  = none            # xor sort
reta      = on              # off
scaling   = on              # off = fixed window
weighting = packet          # flow byte
keys      = 16
channels  = 8 16 20 32 40 64 128
window    = 150000          # used when scaling = off
can_size  = 18750           # packets per channel, scaling = on
\end{lstlisting}

\section{Priebeh simulácie}

Generátor prúdu paketov volá simulátor v~troch momentoch:

\begin{enumerate}
  \item \textbf{Začiatok toku.} Simulátor pridelí toku pamäťový slot. Reťaz RSS vypočíta
        kanál pre každý kľúč a počet kanálov (16~$\times$~7 = 112 hodnôt) a orákulum JSQ
        vyberie svoj kanál.
  \item \textbf{Paket.} Simulátor iba prečíta uložené kanály a pripočíta váhu paketu do
        histogramov. Hash sa pri pakete nepočíta. Keď sa okno daného počtu kanálov naplní,
        vypočítajú sa metriky, zapíše sa riadok výsledkov a vynulujú sa počítadlá tohto
        počtu kanálov.
  \item \textbf{Koniec toku.} Slot sa uvoľní pre ďalší tok, takže pamäť závisí od počtu
        súčasne živých tokov, nie od ich celkového počtu.
\end{enumerate}

\section{Orákulum JSQ}
\label{sec:jsq}

Orákulum \emph{Join the Shortest Queue} (JSQ) slúži ako \textbf{dolná hranica}: ukazuje, čo
sa dá dosiahnuť pri zachovaní pravidla ,,jeden tok = jeden kanál``, ak by jednotka poznala
veľkosť a trvanie toku už pri jeho začiatku. V~hardvéri sa takto realizovať nedá.

Nový tok ide do kanála s~najmenšou očakávanou záťažou na konci aktuálneho okna:
\begin{equation}
  \mathit{ch} = \arg\min_{i} \left( h_i + t_\mathit{left} \cdot r_i \right) ,
\end{equation}
kde $h_i$ je záťaž doručená v~tomto okne, $r_i$ súčet rýchlostí živých tokov na kanáli a
$t_\mathit{left}$ odhad času do konca okna (chýbajúce pakety delené rýchlosťou všetkých živých
tokov). Toky kratšie než 1\,ms sa do rýchlostí nezapočítavajú.

Skúšali sme štyri verzie (syntetické dáta, váha \texttt{packet}, porovnanie s~Toeplitzom):

\begin{description}
  \item[A] kumulatívna záťaž od začiatku: horšie než hash o~4--10\,\%,
  \item[B] záťaž aktuálneho okna: $-19$\,\% pri 8 kanáloch, $+1$\,\% pri 128,
  \item[C] najmenej aktívnych tokov: približne ako hash,
  \item[D] doručená záťaž a očakávaný zvyšok (vzorec vyššie): $-23$\,\% pri 8 kanáloch,
        $-5$\,\% pri 128. Túto verziu sme ponechali.
\end{description}

\section{Overenie správnosti}

Pri prepise simulátora sa výsledky porovnávali s~pôvodnou implementáciou:

\begin{itemize}
  \item na malom testovacom datasete všetkých 180 kombinácií (5 hashov, 2 tuple,
        3 symetrie, RETA zap./vyp., 3 váhy), spolu 108\,360 riadkov,
  \item na celom \texttt{odid1} 212\,660 riadkov (15\,190 okien, 7 počtov kanálov, Toeplitz a JSQ).
\end{itemize}

Najväčší relatívny rozdiel bol menší než $10^{-5}$ a pochádza zo zaokrúhlenia uložených
hodnôt pri priemerovaní cez kľúče. Hodnoty JSQ na \texttt{odid1} sa zhodovali presne.

\TODO{Doplniť čas behu a spotrebu pamäte na \texttt{odid1}.}
